% Source: https://artofproblemsolving.com/community/c3h2008151_arithmetic_problems_from_cut_the_knot
% Stable identity: aops-topic-2008151-q01
Cutting a triangle into $100$ congruent small triangles by lines parallel to the sides means each side is divided into $n$ equal parts with $n^2=100$, so $n=10$. Rows from the top contain $1,3,5,\dots,19$ cells. In each row the cells alternate between \emph{upward} triangles (pointing the same way as the board) and \emph{downward} ones, and the row with $2k-1$ cells has $k$ upward and $k-1$ downward cells. Hence there are
\[
1+2+\cdots+10=55 \text{ upward cells and } 0+1+\cdots+9=45 \text{ downward cells.}
\]
(The figure shows the pattern on a smaller board with $n=4$.)

\begin{asy}
usepackage("amsmath");
size(6cm);
defaultpen(fontsize(10pt));
real h=sqrt(3)/2;
int n=4;
for(int i=0;i<n;++i) for(int j=0;j<n-i;++j){
  pair p=(j+i/2,i*h);
  filldraw(p--p+(1,0)--p+(0.5,h)--cycle,white,black);
  if(j<n-i-1) filldraw(p+(1,0)--p+(0.5,h)--p+(1.5,h)--cycle,gray(0.6),black);
}
label("upward (white)",(4.6,2.6),E);label("downward (gray)",(4.6,2.1),E);
filldraw((4.25,2.5)--(4.55,2.5)--(4.4,2.75)--cycle,white,black);
filldraw((4.25,2.2)--(4.55,2.2)--(4.4,1.95)--cycle,gray(0.6),black);
label("\mbox{}",(min(currentpicture,user=true)+max(currentpicture,user=true))/2,black); // keeps the text after the figure black when LaTeX typesets the labels inline
\end{asy}

Two neighbors share a side, and every side of an upward cell inside the board is a side of a downward cell, and vice versa. So each hop takes a grasshopper from an upward cell to a downward cell or from a downward cell to an upward cell. After $9$ hops (an odd number), each grasshopper is on a cell of the opposite type from its starting cell.

In particular, the grasshoppers now sitting on upward cells are exactly those that started on downward cells, and there are only $45$ of them. Thus at most $45$ of the $55$ upward cells are occupied, so at least $55-45=10$ upward cells are empty. $\blacksquare$
